% Chapter 12: Eleven Minutes
% Scene function: Post-turn haunting. The work continues at full
% quality. But Tert waits eleven minutes instead of four. The shift
% report calls it thermal conditions. Tert knows the difference.

\chapter{Eleven Minutes}

I go back to the house on Monday.

The work does not change. I want to be clear about this. The work does not change because the work is the work, and what happened on Sunday night does not alter the fact that I have an assignment with a timeline and an escalation schedule and a set of professional standards I have spent my career developing. The house still needs to do what I have taught it to do. Gordon still needs to feel what the Phase schedule says he needs to feel.

I produce two sound events on Monday night, both at the bass register, both timed to his first sleep cycle. They are technically clean. On Tuesday I advance the thermal gradient in the hallway by another half-degree, targeting the space between the bedroom and the bathroom. On Wednesday I introduce a new element: a vibration in the attic ladder mechanism that produces a sound like a footstep on the landing, once, at 2:14 a.m., the time of night when the brain is most susceptible to interpreting ambiguous sounds as human presence.

This is good work. This is the work I am here to do.

On Wednesday night, after the attic sound, Gordon wakes. I hear him sit up in bed. I hear the silence that follows, which is the silence of a person listening, deciding whether what they heard was real.

The schedule says: follow the attic sound with a second event within four minutes. This is the escalation principle. You do not let the subject settle. You layer. The second event confirms the first and collapses the possibility that it was imagined.

I wait eleven minutes.

I tell myself this is because the thermal pocket in the hallway has not reached optimal differential. This is true. It is also true that I could have used a different second event, one that did not depend on the thermal pocket, and I chose not to.

Gordon lies back down. He does not go back to sleep for a long time, but he lies down, and the quality of his breathing changes from listening to waiting to something that is not quite rest but is closer to rest than listening was. Outside, through the bedroom window, a car passes on the street. Its headlights move across the ceiling. Gordon watches them.

I file the shift report on Thursday morning. Under activities I write: \emph{Two sound events Mon. Thermal advancement Tues. Attic mechanism Wed, single event, 2:14. Follow-up delayed due to thermal conditions. No contact.}

``Delayed due to thermal conditions'' is not a lie. It is a true statement. I could have used a different follow-up, one that did not depend on the thermal pocket. I chose not to.
